\documentclass[10pt,a4paper,ngerman]{scrartcl}
\usepackage[utf8]{inputenc}
\usepackage[ngerman]{babel}
\usepackage[T1]{fontenc}
\usepackage{amsmath}
\usepackage{amsfonts}
\usepackage{amssymb}
\usepackage{makeidx}
\usepackage{graphicx}
\usepackage{epstopdf}
\usepackage[onehalfspacing]{setspace}
\usepackage{array}
\usepackage{upgreek}
\usepackage{float}
\usepackage{csquotes}
\usepackage{chemformula}
\usepackage{chemgreek}
\usepackage{chemmacros}
\chemsetup{modules=all}
\usepackage{tikz}
\usepackage{siunitx}
\sisetup{locale = DE,per-mode = symbol}
\usepackage{esvect}
\usepackage{geometry}
\usepackage{pdfpages}
\usepackage[font=small]{caption}
%\usepackage[version=4]{mhchem}
\usepackage[backend=biber, style=chem-angew]{biblatex}
\addbibresource{Polyaddition_Polykond.bib}
\usepackage{tabularx, booktabs, multirow}
\usepackage[pdfborder={0 0 0}]{hyperref}
\usepackage{scrlayer-scrpage}%Header für die KOMA-script -Klasse
%\usepackage{hyperref}
\usepackage{pgfplots}
\usepackage{textgreek}


\newcommand{\tildenu}{{\fontencoding{LGR}\selectfont\accperispomeni\textnu}}
\captionsetup{format=plain}
\chemsetup[phases]{pos=sub}
\chemsetup[redox]{pos=top}
\chemsetup[redox]{align=center}
%\chemsetup[redox]{explicit-sign=true}
%\chemsetup[redox]{roman=false}
\newcommand{\celsius}{^{\circ}\mathrm{C}}
\parindent0pt
\sloppy
\DeclareChemPhase{\s}{s}
\DeclareChemPhase{\l}{l}
\DeclareChemPhase{\g}{g}
\renewcaptionname{ngerman}{\figurename}{Abb.}
\renewcaptionname{ngerman}{\tablename}{Tab.}
\geometry{bottom=100pt} \geometry{top=80pt}
\newcommand{\m}{\mathrm{m}}
\newcommand{\K}{\mathrm{K}}
\newcommand{\J}{\mathrm{J}}
\newcommand{\gr}{\mathrm{g}}
\newcommand{\kg}{\mathrm{kg}}
\newcommand{\mol}{\mathrm{mol}}
\newcommand{\Pa}{\mathrm{Pa}}
\newcommand{\ad}{\mathrm{ad}}
\newcommand{\de}{\mathrm{de}}
\newcommand{\gmol}{\mathrm{\frac{g}{mol}}}
\newcommand{\JKmol}{\mathrm{\frac{J}{K \cdot mol}}}
\newcommand{\kgmss}{\mathrm{\frac{kg}{m \cdot s^2}}}
\newcommand{\loge}{\mathrm{ln}}


\begin{document}
\begin{titlepage}
	\vspace*{2,5cm}
	\begin{center}
		\begin{LARGE}
			\textbf{Polymer Chemie}\\
			\textbf{Viscometry}\\
		\end{LARGE}
		\vspace{1cm}
		\textbf{Universität Stuttgart}\\
		\vspace*{1,5cm}
		\begin{tabular}{lp{7,5cm}l}
			Author 1:     &Laurens Viehoff\\
            Author 2:     &Felix Roth\\   
			E-Mail 1:      &st183688@stud.uni-stuttgart.de\\
            E-Mail 2:      &st188157@stud.uni-stuttgart.de  \\
			Student-ID 1: & 3676761   \\ 
            Student-ID 2: & 3711192 \\ \\
			Assistentin:  &   \\ \\ 
		\end{tabular}
	\end{center}

    
\end{titlepage}
\thispagestyle{empty}
\renewcommand{\thepage}{\arabic{page}}
\newpage
\tableofcontents
\setcounter{page}{1}
\renewcommand{\thepage}{\arabic{page}}
\newpage



\section{Theory}	

	\subsection{Viscosity}

		Viscosity in polymers is influenced by 5 key factors.
		The first being the chain lenght of the Polymer, the longer the chains the more viscouse the fluid gets.
		This is due to the Intramolecular forces.
		The second influence is the chemical structure of the polymer, different functional groups mean different intermolecular forces.
		The third one is the influence of the solvent. Because different solvents have different properties.
		The fourth is the influence of the concentration. This means with different concentration the viscosity changes.
		The last but not least influence is the temperature. Because the viscosity changes with different temperatures because the movement of the molecules changes.\\

	\subsection{From viscosity to $M_\eta$}

		The viscosity is defined after Hagen-Poiseuille by following equation:

			\begin{equation}
				\eta = \frac{\pi\cdot r^4 \cdot t \cdot \Delta p}{8\cdot l \cdot V}
				\label{viscosity}
			\end{equation}

		In which the $\Delta p$ is defindet as:
			
			\begin{equation}
				\Delta p = \rho \cdot g \cdot h
			\end{equation}

		In these equations $r$ is the capillaryradius, $t$ is the time the polymer needs to flow through the capillary,
		$\Delta p$ is the pressure difference in the capillary, $l$ is the lenght of the capillary, $V$ is the Volume, 
		$\rho$ is the density of the solution and $h$ is the height of the measured part.

		After analyzing different polymers four different types of Viscosities could be determined.

			\begin{eqnarray}
				\eta_\text{rel} =& \frac{\eta}{\eta_0}\\
				\eta_\text{sp} =& \eta_\text{rel} - 1^\\ 
				\eta_\text{red} =& \frac{\eta_\text{sp}}{c}\\ 
				\eta_\text{} =& \frac{\ln \eta_\text{rel}}{c}
			\end{eqnarray}

		 
		If equation \ref{viscosity} is taken into account $\eta_\text{sp}$ is changed to:
		
			\begin{equation}
				\eta_\text{sp} = \frac{t\rho - t_0\rho_0}{t_0\rho_0}
			\end{equation}

		And because of the low concentration the density of the solution is almost equal to the density of the solvent.
		Therefore the equation can be simplyfied.
			
			\begin{equation}
				\eta_\text{sp} \approx \frac{t-t_0}{t_0}
			\end{equation}

	\newpage

		To get the viscousmiddeled molecular Volume the Staudinger-index is used, and it is connected via the 
		Marc-Howink-Equation.
		
			\begin{equation}
				[\eta] = K \cdot \bar{M_\eta^a}
			\end{equation}

		This equation is used in combination with the equatoin for the viscositynumber.

			\begin{equation}
				[\eta] = \lim_{c\rightarrow0} (\eta_\text{red}) =\lim_{c\rightarrow0} \left( \frac{\eta_\text{sp}}{c}\right) 
			\end{equation}
		

		So by extrapolating the results $\eta_\text{red}$ can be determined and via this $[\eta]$ and from that $\bar{M_\eta}$ can be determined.
		This is also used in our case in the experiment.


\section{Experimental}

		The three different polymer solutions were already prepared and ready to use. 
		One pure Toluene solution and thre polymer solutions with different concentrations.
		all of these four solutions were analyzed with a Viscosimeter.

\section{Results}

		


\section{Questions}

		\subsection{Question 1}

			In total it can be stated, all three of the given polymers have different properties.
			Polystyrene is a neutral polymer and behaves like a statisticall tangle in a good solvent, the Staudinger-index
			only is a referencevalue. The Polystyrene-lithium is a "living" anionic polymer with Lithium as a counter ion.
			This Lithium-ion forms a polar Li-C-Bond in polar solvents, these are also called Units.
			These Units have a direct influence on the Staudinger-Index, so the Index gets higher.
			In the last case of the sulfonated Polystyrene is a bit mor komplex. 
			PSS-Na is a Polyelectrolyte, therefore the \ch{SO3-} is distributed over the chain. 
			A Uniqueness is if this Polymer is put in demineralized water the negative charges repel each other.
			Therefore the polymer swells and the tangle gets bigger, the Staudinger index gets very high and is the highest from these three examples.


		\subsection{Question 2}

			At a certain concentration point things like the intramolecular forces, the hydrodynamic interactions and the entanglements cannot longer be ignored.
			This is because at low concentrations the dilution is assumed to be infinite, therefore the polymer molecule is "alone" and has no ineractions with other molecules.


		\subsection{Question 3}

			The biggest difference between these two viscometers is the third vent pipe at the Ubeholde viscosimeters.
			This has the effect, that for the Ubeholde-viscometer the  of the sample is almost irrelevant.
			It also helps to get a constant hydrostatic pressure. 
			Therefore the use cases vary alot between those two. The Ostwald-voscometer is used for individual measurements and the Ubeholde-viscometer is used for dialation measurements.


		\subsection{Question 4}

			The Staudinger-Index of the it-PMMA is higher because of the regulary orderd groups the molecule is sterically hindered. 
			The at-PMMA has a random order therefore there is less sterically hindrence.
			Another factor is the solvent because the it-PMMA is already more "swallen" the effect of the solvant is even stronger. 


		

\printbibliography

\end{document}